\chapter{MARCO TEÓRICO}

Este capítulo desarrolla los fundamentos conceptuales y formales sobre los que se
construye DARL. La presentación sigue el orden lógico del framework: primero se
describe el dominio de aplicación de la tesis (datos tabulares), luego se formaliza
el objeto de estudio (el \textit{pipeline} de dos etapas), después se introduce el
fenómeno que perturba su funcionamiento (\textit{distribution shift}), posteriormente
se presentan las señales que permiten observarlo (métricas de monitoreo), y
finalmente se desarrolla el formalismo de decisión y el algoritmo de aprendizaje
que constituyen el núcleo de la contribución (POMDP y PPO).

% ------------------------------------------------------------
\section{Datos tabulares en aprendizaje automático}
% ------------------------------------------------------------

Los datos tabulares constituyen una de las formas más comunes de representación de
información estructurada en aplicaciones reales de aprendizaje automático. En este
tipo de datos, cada observación se organiza como una fila de una tabla y cada
variable o atributo se representa como una columna. A diferencia de dominios como
visión por computadora o procesamiento de lenguaje natural, donde las entradas
suelen tener una estructura homogénea, los datos tabulares combinan variables
numéricas, categóricas, ordinales, binarias y valores faltantes dentro de una misma
representación.

Formalmente, un conjunto de datos tabular supervisado puede representarse como:

\begin{equation}
    \mathcal{D} = \{(x_i, y_i)\}_{i=1}^{n}, \qquad
    x_i = (x_{i1}, x_{i2}, \ldots, x_{id}) \in \mathcal{X}.
    \label{eq:tabular_dataset}
\end{equation}

\noindent\textbf{Donde:}
\begin{itemize}
    \item $\mathcal{D}$ representa el conjunto de datos disponible para entrenamiento, validación o evaluación.
    \item $n$ representa el número total de observaciones o filas de la tabla.
    \item $x_i$ representa la observación tabular número $i$.
    \item $y_i$ representa la etiqueta o variable objetivo asociada a la observación $x_i$.
    \item $x_{ij}$ representa el valor de la variable o columna $j$ para la observación $i$.
    \item $d$ representa el número de variables originales en la tabla.
    \item $\mathcal{X}$ representa el espacio de entrada formado por las observaciones tabulares crudas.
\end{itemize}

Conceptualmente, esta notación muestra que una tabla no es solo una matriz de valores,
sino una colección de observaciones con atributos heterogéneos y una variable objetivo.
Esta heterogeneidad es relevante porque distintas columnas requieren tratamientos
diferentes antes de ser usadas por un modelo predictivo. Por ejemplo, las variables
numéricas pueden requerir escalamiento o transformación de cuantiles, mientras que las
variables categóricas pueden requerir codificación, agrupamiento de categorías raras o
manejo explícito de valores desconocidos.

En consecuencia, el desempeño de un modelo tabular no depende únicamente del
algoritmo predictivo, sino también de las decisiones de preprocesamiento aplicadas antes
del entrenamiento. Esta característica distingue a los sistemas tabulares de otros dominios
donde las representaciones suelen aprenderse de manera más integrada. En problemas
tabulares, la etapa de \textit{feature engineering} puede modificar la escala, la dimensión y
la semántica de las variables que recibe el modelo.

TableShift es especialmente relevante para esta investigación porque fue propuesto como
un \textit{benchmark} para estudiar \textit{distribution shift} en tareas de clasificación
binaria con datos tabulares reales \cite{gardner2023tableshift}. Este \textit{benchmark}
reúne tareas de distintos dominios, como salud, finanzas, educación, política pública y
participación cívica, y permite evaluar cómo cambia el desempeño de los modelos cuando
la distribución de prueba difiere de la distribución de entrenamiento. Por ello, TableShift
sirve como referencia conceptual y experimental para estudiar la robustez de modelos
tabulares ante cambios de distribución.

Para DARL, la importancia de los datos tabulares no se limita al tipo de datos utilizado,
sino que define el problema de mantenimiento del \textit{pipeline}. Si cambian las
frecuencias de categorías, los rangos numéricos, la proporción de valores faltantes o las
relaciones entre variables, los parámetros aprendidos por el preprocesamiento pueden
quedar desactualizados. Esto justifica separar explícitamente la etapa de transformación
de características y la etapa predictiva, como se formaliza en la siguiente sección.

% ------------------------------------------------------------
\section{Pipelines de ML de dos etapas sobre datos tabulares}
% ------------------------------------------------------------
 
Un \textit{pipeline} de aprendizaje automático sobre datos tabulares es un sistema
compuesto por una secuencia de operaciones que transforma observaciones crudas en
predicciones. En esta investigación se considera un \textit{pipeline} de dos etapas:
una etapa de \textit{feature engineering} o preprocesamiento, encargada de convertir
los datos tabulares originales en una representación utilizable por el modelo, y una
etapa de modelo predictivo, encargada de aprender la relación entre dicha representación
y la variable objetivo.

Esta separación es importante porque ambos componentes pueden verse afectados de
manera distinta ante cambios de distribución. Un cambio en la distribución marginal de
las variables de entrada puede volver obsoletas las estadísticas aprendidas por el
preprocesamiento, mientras que un cambio en la relación entre entrada y salida puede
volver obsoletos los parámetros del modelo predictivo. Por ello, DARL analiza el
\textit{pipeline} como una composición de dos funciones con parámetros diferentes.
 
La primera etapa, denominada \textit{feature engineering} o preprocesamiento, puede
formalizarse mediante una función parametrizada:

\begin{equation}
    f_\phi: \mathcal{X} \rightarrow \mathcal{Z}.
    \label{eq:feature_engineering_mapping}
\end{equation}

\noindent\textbf{Donde:}
\begin{itemize}
    \item $\mathcal{X} \subseteq \mathbb{R}^{d}$ corresponde al espacio de entrada original, compuesto por los datos tabulares crudos. Se define como un vector de $d$ variables o columnas.
    \item $\mathcal{Z} \subseteq \mathbb{R}^{k}$ representa el espacio transformado o representación intermedia utilizada por el modelo predictivo con $k$ variables, que pueden diferir de las $d$ variables originales.
    \item $f_\phi$ representa la función de preprocesamiento o ingeniería de características.
    \item $\phi$ representa los parámetros asociados a dicha función, tales como medias, desviaciones estándar, codificaciones categóricas, cuantiles o reglas de imputación aprendidas durante el entrenamiento.
\end{itemize}

Conceptualmente, esta notación muestra que la etapa de \textit{feature\ engineering} no es solo una transformación fija, sino un proceso parametrizado que puede modificar la dimensión y la escala de los datos. Esto es especialmente importante en datos tabulares, donde las variables numéricas, categóricas y ordinales requieren tratamientos diferentes antes de ser usadas por un modelo predictivo. Así, herramientas como escaladores, imputadores o codificadores aprenden valores a partir del entrenamiento para transformar nuevos datos; por ello, ante cambios en la distribución de entrada, los parámetros ($\phi$) pueden volverse obsoletos y requerir actualización.
 
La segunda etapa, denominada modelo predictivo, se define como una función que toma
la representación intermedia y produce una salida:

\begin{equation}
    g_\theta: \mathcal{Z} \rightarrow \mathcal{Y}.
    \label{eq:model_mapping}
\end{equation}

\noindent\textbf{Donde:}
\begin{itemize}
    \item $\mathcal{Z} \subseteq \mathbb{R}^{k}$ representa el espacio de características ya transformadas por la etapa de preprocesamiento.
    \item $\mathcal{Y} \subseteq \mathbb{R}$ representa el espacio de salida, como las clases resultantes.
    \item $g_\theta$ representa el modelo predictivo entrenado sobre las características transformadas.
    \item $\theta$ representa los parámetros aprendidos del modelo, como pesos, umbrales, estructuras de árboles o hiperparámetros ajustados durante el entrenamiento.
\end{itemize}

Conceptualmente, esta fórmula separa claramente el rol del preprocesamiento y el rol
del modelo. Mientras $f_\phi$ transforma los datos, $g_\theta$ aprende la relación entre
la representación transformada y la variable objetivo. En esta investigación, el modelo
predictivo principal es XGBoost \cite{chen2016xgboost}, con regresión logística como
línea base.
 
La función de predicción del \textit{pipeline} completo se obtiene componiendo ambas
etapas:

\begin{equation}
    h = g_\theta \circ f_\phi, \qquad
    \hat{y} = h(x) = g_\theta\!\left(f_\phi(x)\right).
    \label{eq:pipeline_composition}
\end{equation}

\noindent\textbf{Donde:}
\begin{itemize}
    \item $h$ representa la función completa del \textit{pipeline} de ML.
    \item $\circ$ representa la composición de funciones, es decir, aplicar primero $f_\phi$ y luego $g_\theta$.
    \item $x$ representa una observación tabular cruda.
    \item $f_\phi(x)$ representa la versión transformada de la observación $x$.
    \item $g_\theta(f_\phi(x))$ representa la predicción producida por el modelo a partir de la observación transformada.
    \item $\hat{y}$ representa la salida predicha por el \textit{pipeline}.
\end{itemize}

Esta composición expresa que la predicción final depende tanto de la
calidad de la transformación como de la calidad del modelo. Si el rendimiento de $h$ se
degrada, el origen del problema puede estar en $f_\phi$, en $g_\theta$ o en ambos. Esta
separación es la base del problema de actualización selectiva que DARL busca resolver:
identificar qué componente debe actualizarse bajo distintos tipos de drift.
 
% ------------------------------------------------------------
\section{Distribution shift: taxonomía y formalización}
% ------------------------------------------------------------
 
Sea $P_{\text{tr}}(X,Y)$ la distribución conjunta de entrada y salida durante el
entrenamiento, y sea $P_{\text{te}}(X,Y)$ la distribución conjunta durante el despliegue.
De forma general, existe \textit{distribution shift} cuando:
 
\begin{equation}
    P_{\text{tr}}(X,Y) \neq P_{\text{te}}(X,Y).
    \label{eq:distribution_shift}
\end{equation}
 
\noindent\textbf{Donde:}
\begin{itemize}
    \item $X$ representa las variables de entrada o características del problema.
    \item $Y$ representa la variable objetivo o etiqueta que se desea predecir.
    \item $P_{\text{tr}}(X,Y)$ representa la distribución conjunta observada durante el entrenamiento.
    \item $P_{\text{te}}(X,Y)$ representa la distribución conjunta observada durante el despliegue o evaluación.
\end{itemize}

Conceptualmente, esta fórmula indica que el modelo fue entrenado bajo una distribución
estadística distinta a la que enfrenta durante producción. Esta diferencia puede hacer que
un modelo con buen desempeño durante entrenamiento pierda precisión en despliegue,
incluso si su arquitectura y sus parámetros no han cambiado \cite{gama2014survey,cai2023disde}.

La distribución conjunta puede factorizarse en una distribución marginal de entrada y
una distribución condicional de salida:

\begin{equation}
    P(X,Y) = P(X)\,P(Y \mid X).
    \label{eq:joint_factorization}
\end{equation}

\noindent\textbf{Donde:}
\begin{itemize}
    \item $P(X,Y)$ representa la distribución conjunta de entradas y salidas.
    \item $P(X)$ representa la distribución marginal de las variables de entrada.
    \item $P(Y \mid X)$ representa la distribución condicional de la variable objetivo dado el valor de las entradas.
\end{itemize}

Conceptualmente, esta descomposición permite distinguir dos fuentes principales de
drift. Puede cambiar la distribución de los datos que llegan al sistema, lo que afecta a
$P(X)$, o puede cambiar la relación entre esos datos y la etiqueta, lo que afecta a
$P(Y \mid X)$. Esta distinción es central porque cada tipo de cambio sugiere una acción
de mantenimiento distinta sobre el \textit{pipeline} \cite{cai2023disde}.
 
\subsection{Covariate shift}
 
El \textit{covariate shift} ocurre cuando cambia la distribución marginal de las variables
de entrada, pero se mantiene estable la relación condicional entre entrada y salida:
 
\begin{equation}
    P_{\text{tr}}(X) \neq P_{\text{te}}(X)
    \quad \text{con} \quad
    P_{\text{tr}}(Y \mid X) = P_{\text{te}}(Y \mid X).
    \label{eq:covariate_shift}
\end{equation}
 
\noindent\textbf{Donde:}
\begin{itemize}
    \item $P_{\text{tr}}(X)$ representa la distribución de las variables de entrada durante el entrenamiento.
    \item $P_{\text{te}}(X)$ representa la distribución de las variables de entrada durante el despliegue.
    \item $P_{\text{tr}}(Y \mid X)$ representa la relación entre entrada y salida aprendida en entrenamiento.
    \item $P_{\text{te}}(Y \mid X)$ representa la relación entre entrada y salida vigente durante el despliegue.
    \item La primera desigualdad indica que los datos de entrada cambiaron.
    \item La igualdad condicional indica que la lógica que relaciona $X$ con $Y$ se mantiene estable.
\end{itemize}

Conceptualmente, el \textit{covariate shift} significa que cambian los datos, pero no
necesariamente cambia la regla del problema. En un \textit{pipeline} de dos etapas, este
tipo de drift afecta directamente al preprocesamiento: los parámetros $\phi$, como medias,
desviaciones estándar, cuantiles o frecuencias categóricas, fueron estimados sobre
$P_{\text{tr}}(X)$ y pueden quedar desactualizados cuando la distribución de entrada cambia.
Por ello, actualizar $f_\phi$ puede ser suficiente si $g_\theta$ todavía representa
adecuadamente la relación $P(Y \mid X)$ \cite{gama2014survey,cai2023disde}.
 
\subsection{Concept drift}
 
El \textit{concept drift} ocurre cuando cambia la relación condicional entre las variables
de entrada y la variable objetivo, incluso si la distribución de entrada permanece estable:
 
\begin{equation}
    P_{\text{tr}}(Y \mid X) \neq P_{\text{te}}(Y \mid X)
    \quad \text{con} \quad
    P_{\text{tr}}(X) = P_{\text{te}}(X).
    \label{eq:concept_drift}
\end{equation}
 
\noindent\textbf{Donde:}
\begin{itemize}
    \item $P_{\text{tr}}(Y \mid X)$ representa la relación entrada-salida observada durante el entrenamiento.
    \item $P_{\text{te}}(Y \mid X)$ representa la relación entrada-salida vigente durante el despliegue.
    \item $P_{\text{tr}}(X)$ representa la distribución de entradas en entrenamiento.
    \item $P_{\text{te}}(X)$ representa la distribución de entradas en despliegue.
    \item La desigualdad condicional indica que cambió la función objetivo del problema.
    \item La igualdad marginal indica que las entradas pueden verse estadísticamente similares aunque el significado predictivo haya cambiado.
\end{itemize}

Conceptualmente, el \textit{concept drift} implica que la lógica del problema cambió. En
el \textit{pipeline} de dos etapas, esto afecta principalmente al modelo predictivo:
los parámetros $\theta$ aprendidos sobre $P_{\text{tr}}(Y \mid X)$ ya no aproximan
correctamente la relación vigente $P_{\text{te}}(Y \mid X)$. En este caso, actualizar
solo $f_\phi$ suele ser insuficiente, porque el problema no está en la escala o codificación
de los datos, sino en la relación que el modelo debe aprender \cite{cai2023disde}.
 
\subsection{Drift combinado y taxonomía de severidad}
 
En la práctica, ambos tipos de drift pueden coexistir \cite{gama2014survey}. Sea cual
sea el tipo, el drift puede clasificarse también según su patrón temporal:
 
\begin{itemize}
    \item Drift abrupto: la distribución cambia de manera repentina en un único instante $t^*$.
    \item Drift gradual: la distribución transiciona suavemente entre $P_{\text{tr}}$ y $P_{\text{te}}$ a lo largo de un intervalo de tiempo.
    \item Drift incremental: la distribución se aleja monotónicamente de $P_{\text{tr}}$ sin retornar.
    \item Drift recurrente: la distribución oscila entre estados conocidos a lo largo del tiempo.
\end{itemize}
 
Esta investigación trabaja con drift abrupto de tipo y severidad controlados para aislar
el efecto de cada condición experimental \cite{gardner2023tableshift,shakhovska2025severity}.
 
% ------------------------------------------------------------
\section{Métricas de monitoreo de drift}
% ------------------------------------------------------------
 
El módulo de monitoreo de DARL construye, en cada instante de decisión, un vector de
observación a partir de métricas estadísticas calculadas sobre ventanas de datos:

\begin{equation}
    \mathbf{o}_t \in \mathbb{R}^{m}.
    \label{eq:observation_vector_generic}
\end{equation}

\noindent\textbf{Donde:}
\begin{itemize}
    \item $\mathbf{o}_t$ representa el vector de observación disponible para el agente en el instante $t$.
    \item $m$ representa el número total de métricas o señales incluidas en la observación.
    \item $\mathbb{R}^{m}$ indica que la observación se representa como un vector real de dimensión $m$.
\end{itemize}

Conceptualmente, el agente DARL no observa directamente el estado real de la distribución,
sino un resumen estadístico construido a partir de los datos disponibles. Este vector reúne
señales de cambio distribucional, severidad y posible degradación, y actúa como entrada del
POMDP formalizado en la Sección~\ref{sec:pomdp}.
 
\subsection{Índice de estabilidad poblacional (PSI)}
 
El \textit{Population Stability Index} (PSI) cuantifica el cambio en la distribución
marginal de una variable continua entre una ventana de referencia y una ventana actual
\cite{shakhovska2025severity}. Para una variable particionada en $n$ intervalos, el PSI
se define como:
 
\begin{equation}
    \text{PSI} = \sum_{i=1}^{n} (p_i - q_i)\,\ln\!\left(\frac{p_i}{q_i}\right).
    \label{eq:psi}
\end{equation}
 
\noindent\textbf{Donde:}
\begin{itemize}
    \item $\text{PSI}$ representa el índice de estabilidad poblacional de una variable.
    \item $n$ representa el número de intervalos o bins usados para comparar las distribuciones.
    \item $i$ representa el índice de cada intervalo.
    \item $p_i$ representa la proporción de observaciones de entrenamiento que caen en el intervalo $i$.
    \item $q_i$ representa la proporción de observaciones actuales o de producción que caen en el intervalo $i$.
    \item $\ln(\cdot)$ representa el logaritmo natural.
\end{itemize}

Conceptualmente, el PSI compara cómo se reparte una variable entre intervalos antes y
después del despliegue. Si las proporciones $p_i$ y $q_i$ son similares en todos los
intervalos, el valor del PSI será bajo. Si una variable empieza a concentrarse en rangos
distintos a los observados durante entrenamiento, el PSI aumenta. En DARL se calcula un
valor de PSI por cada \textit{feature}, lo que permite identificar qué variables presentan
mayor desplazamiento. Como regla práctica, un $\text{PSI} < 0.1$ indica estabilidad,
valores entre $0.1$ y $0.2$ sugieren cambio moderado, y $\text{PSI} > 0.2$ indica un
cambio significativo.
 
\subsection{Test de Kolmogorov--Smirnov (KS)}
 
El test de Kolmogorov--Smirnov es una prueba no paramétrica de igualdad de
distribuciones que compara las funciones de distribución acumulada empíricas
$\hat{F}_{\text{tr}}$ y $\hat{F}_{\text{te}}$ \cite{shakhovska2025severity}. El estadístico
se define como:
 
\begin{equation}
    D_{\text{KS}} = \sup_{x \in \mathbb{R}}
    \left|\hat{F}_{\text{tr}}(x) - \hat{F}_{\text{te}}(x)\right|.
    \label{eq:ks}
\end{equation}
 
\noindent\textbf{Donde:}
\begin{itemize}
    \item $D_{\text{KS}}$ representa el estadístico del test de Kolmogorov--Smirnov.
    \item $\sup$ representa el supremo, es decir, la mayor diferencia observada entre ambas funciones acumuladas.
    \item $x$ representa un posible valor de la variable analizada.
    \item $\mathbb{R}$ representa el conjunto de valores reales sobre los que se evalúa la variable.
    \item $\hat{F}_{\text{tr}}(x)$ representa la función de distribución acumulada empírica de la ventana de entrenamiento.
    \item $\hat{F}_{\text{te}}(x)$ representa la función de distribución acumulada empírica de la ventana actual o de despliegue.
    \item $|\cdot|$ representa el valor absoluto de la diferencia entre ambas funciones acumuladas.
\end{itemize}

Conceptualmente, el estadístico KS mide la mayor separación entre dos distribuciones
acumuladas. A diferencia del PSI, no requiere discretizar la variable en intervalos y puede
ser sensible a diferencias en cualquier región de la distribución. En DARL se calcula
$D_{\text{KS}}$ por cada \textit{feature} numérico para detectar señales de \textit{covariate shift}
de forma individual.
 
\subsection{Classifier Two-Sample Test (C2ST)}
 
El \textit{Classifier Two-Sample Test} (C2ST) es un test de dos muestras basado en
clasificación que detecta diferencias distribucionales entre el conjunto de entrenamiento
y la ventana actual \cite{cai2023disde,singh2025shift}. El procedimiento construye un
problema binario en el que un clasificador intenta distinguir si una observación proviene
del periodo de entrenamiento o del periodo actual. El estadístico de test se define como:
 
\begin{equation}
    \text{C2ST} = \text{AUC-ROC}(c).
    \label{eq:c2st}
\end{equation}
 
\noindent\textbf{Donde:}
\begin{itemize}
    \item $\text{C2ST}$ representa el estadístico del test basado en clasificación.
    \item $c$ representa el clasificador entrenado para distinguir entre datos de entrenamiento y datos actuales.
    \item $\text{AUC-ROC}(c)$ representa el área bajo la curva ROC obtenida por el clasificador $c$.
\end{itemize}

Conceptualmente, si el clasificador $c$ no puede distinguir de qué periodo provienen las
observaciones, entonces las distribuciones son similares y el valor de C2ST se acerca a
$0.5$. En cambio, si el clasificador separa claramente ambos periodos, el valor se aproxima
a $1$, lo que indica una diferencia distribucional relevante. En DARL, el C2ST puede usarse
como una señal global de separación entre periodos; cuando se aplica únicamente sobre
$X$, detecta cambios en la distribución de entrada, y cuando se dispone de etiquetas
retrasadas o señales derivadas de desempeño, puede contribuir al diagnóstico de cambios
asociados a la relación $P(Y \mid X)$.
 
\subsection{Severity score combinado}
 
Las métricas anteriores se agregan en un \textit{severity score} que cuantifica la magnitud
global del drift observado en el instante $t$ \cite{shakhovska2025severity}:
 
\begin{equation}
    S_t = \alpha\,\overline{\text{PSI}}_t
    + \beta\,\overline{D}_{\text{KS},t}
    + \gamma\,\text{C2ST}_t,
    \qquad
    \alpha + \beta + \gamma = 1.
    \label{eq:severity_score}
\end{equation}
 
\noindent\textbf{Donde:}
\begin{itemize}
    \item $S_t$ representa el puntaje combinado de severidad del drift en el instante $t$.
    \item $\overline{\text{PSI}}_t$ representa el promedio de los valores PSI por \textit{feature} en el instante $t$.
    \item $\overline{D}_{\text{KS},t}$ representa el promedio de los estadísticos KS por \textit{feature} en el instante $t$.
    \item $\text{C2ST}_t$ representa el valor del test C2ST calculado para la ventana actual.
    \item $\alpha$, $\beta$ y $\gamma$ son pesos que controlan la importancia relativa de cada métrica.
    \item La restricción $\alpha + \beta + \gamma = 1$ garantiza que los pesos formen una combinación convexa.
\end{itemize}

Conceptualmente, el \textit{severity score} resume distintas señales de drift en un único
valor interpretable. Esto permite que el agente no tome decisiones a partir de métricas
aisladamente, sino a partir de una medida agregada de severidad. En DARL, $S_t$ se usa
para categorizar el drift en niveles leve, moderado y severo, los cuales forman parte del
vector de observación del agente.
 
% ------------------------------------------------------------
\section{Aprendizaje por refuerzo}
% ------------------------------------------------------------
 
El aprendizaje por refuerzo (\textit{Reinforcement Learning}, RL) es un paradigma de ML
en el que un agente aprende a tomar decisiones mediante la interacción con un entorno,
maximizando una señal de recompensa acumulada \cite{sutton2018rl}. A diferencia del
aprendizaje supervisado, el agente no recibe ejemplos etiquetados de las acciones correctas,
sino que debe descubrirlas mediante exploración y explotación.
 
\subsection{Proceso de Decisión de Markov (MDP) y fundamentos del RL}
\label{sec:mdp}

El marco matemático estándar del RL es el Proceso de Decisión de Markov (MDP), definido
formalmente como la tupla:

\begin{equation}
    \mathcal{M} = \langle \mathcal{S}, \mathcal{A}, P, R, \gamma \rangle.
    \label{eq:mdp_tuple}
\end{equation}

\noindent\textbf{Donde:}
\begin{itemize}
    \item $\mathcal{M}$ representa el proceso de decisión de Markov.
    \item $\mathcal{S}$ es el espacio de estados posibles del entorno.
    \item $\mathcal{A}$ es el espacio de acciones disponibles para el agente.
    \item $P$ es la función de transición entre estados.
    \item $R$ es la función de recompensa.
    \item $\gamma$ es el factor de descuento de recompensas futuras.
\end{itemize}

Conceptualmente, un MDP modela una situación en la que un agente observa un estado,
toma una acción, recibe una recompensa y pasa a un nuevo estado. Esta estructura permite
formular problemas de decisión secuencial, como el mantenimiento de un \textit{pipeline}
de ML, donde cada acción tiene efectos sobre el rendimiento futuro del sistema
\cite{sutton2018rl}.

La función de transición se expresa como:

\begin{equation}
    P(s' \mid s,a) = \Pr(S_{t+1}=s' \mid S_t=s, A_t=a).
    \label{eq:transition_function}
\end{equation}

\noindent\textbf{Donde:}
\begin{itemize}
    \item $P(s' \mid s,a)$ representa la probabilidad de llegar al estado $s'$ después de tomar la acción $a$ en el estado $s$.
    \item $\Pr(\cdot)$ representa una probabilidad.
    \item $S_t$ representa el estado del entorno en el instante $t$.
    \item $A_t$ representa la acción ejecutada por el agente en el instante $t$.
    \item $S_{t+1}$ representa el estado alcanzado en el siguiente instante.
\end{itemize}

Conceptualmente, esta fórmula indica que las acciones no determinan siempre un único
resultado, sino una distribución de posibles estados futuros. En DARL, una acción como
actualizar el preprocesamiento, actualizar el modelo o reentrenar todo el \textit{pipeline}
puede conducir a distintos niveles de recuperación de rendimiento dependiendo del tipo y
severidad del drift.

La función de recompensa se define como:

\begin{equation}
    R: \mathcal{S} \times \mathcal{A} \rightarrow \mathbb{R}.
    \label{eq:reward_function}
\end{equation}

\noindent\textbf{Donde:}
\begin{itemize}
    \item $R$ representa la función que asigna recompensa a un par estado-acción.
    \item $\mathcal{S}$ representa el espacio de estados.
    \item $\mathcal{A}$ representa el espacio de acciones.
    \item $\mathbb{R}$ representa el conjunto de números reales.
\end{itemize}

Conceptualmente, la recompensa resume qué tan conveniente fue una acción en un estado
determinado. Para DARL, esta recompensa puede combinar recuperación de rendimiento,
costo computacional y penalizaciones por decisiones innecesarias.

Un MDP satisface la propiedad de Markov:

\begin{equation}
    P(s_{t+1} \mid s_t,a_t,\ldots,s_0,a_0)
    = P(s_{t+1} \mid s_t,a_t).
    \label{eq:markov_property}
\end{equation}

\noindent\textbf{Donde:}
\begin{itemize}
    \item $s_t$ representa el estado actual.
    \item $a_t$ representa la acción actual.
    \item $s_{t+1}$ representa el estado futuro inmediato.
    \item $s_0,a_0,\ldots$ representan estados y acciones pasadas.
\end{itemize}

Conceptualmente, la propiedad de Markov establece que el futuro depende del estado
presente y de la acción actual, no de toda la historia previa. En la práctica, esta propiedad
puede no cumplirse perfectamente si el agente no observa toda la información relevante,
lo que motiva la formulación POMDP usada posteriormente.

La política del agente se define como una distribución sobre acciones condicionada al
estado:

\begin{equation}
    \pi(a \mid s) = \Pr(A_t=a \mid S_t=s).
    \label{eq:policy}
\end{equation}

\noindent\textbf{Donde:}
\begin{itemize}
    \item $\pi$ representa la política o regla de decisión del agente.
    \item $\pi(a \mid s)$ representa la probabilidad de elegir la acción $a$ dado el estado $s$.
    \item $A_t$ representa la acción tomada en el instante $t$.
    \item $S_t$ representa el estado observado en el instante $t$.
\end{itemize}

Conceptualmente, la política define el comportamiento del agente. En DARL, la política
aprende cuándo diferir, cuándo actualizar solo el preprocesamiento, cuándo actualizar el
modelo y cuándo reentrenar todo el \textit{pipeline}.

La función de retorno representa la suma descontada de recompensas futuras:

\begin{equation}
    G_t = r_t + \gamma r_{t+1} + \gamma^2 r_{t+2} + \cdots
    = \sum_{k=0}^{\infty} \gamma^k r_{t+k},
    \qquad 0 \leq \gamma \leq 1.
    \label{eq:return}
\end{equation}

\noindent\textbf{Donde:}
\begin{itemize}
    \item $G_t$ representa el retorno acumulado desde el instante $t$.
    \item $r_t$ representa la recompensa recibida en el instante $t$.
    \item $r_{t+k}$ representa la recompensa recibida $k$ pasos después del instante $t$.
    \item $\gamma$ representa el factor de descuento.
    \item $k$ representa el número de pasos futuros considerados en la suma.
\end{itemize}

Conceptualmente, el retorno permite evaluar una decisión no solo por su beneficio inmediato,
sino también por sus consecuencias futuras. En DARL, una acción con bajo costo inmediato
podría ser mala si permite que el modelo siga degradándose, mientras que una acción costosa
podría ser conveniente si recupera el rendimiento de forma sostenida.

El objetivo del agente es aprender una política óptima:

\begin{equation}
    \pi^* = \arg\max_{\pi} \; \mathbb{E}_{\pi}\!\left[G_t\right].
    \label{eq:objetivo_rl}
\end{equation}

\noindent\textbf{Donde:}
\begin{itemize}
    \item $\pi^*$ representa la política óptima.
    \item $\arg\max_{\pi}$ representa la política que maximiza el objetivo.
    \item $\mathbb{E}_{\pi}[\cdot]$ representa la esperanza bajo trayectorias generadas por la política $\pi$.
    \item $G_t$ representa el retorno acumulado desde el instante $t$.
\end{itemize}

Conceptualmente, esta fórmula indica que el agente no busca maximizar una recompensa
inmediata aislada, sino el rendimiento esperado acumulado. En DARL, esto se traduce en
aprender una estrategia de mantenimiento que equilibre recuperación de desempeño y costo
de actualización a lo largo del tiempo \cite{sutton2018rl}.

La función de valor de estado mide la recompensa acumulada esperada al encontrarse en
un estado $s$ y seguir la política $\pi$:

\begin{equation}
    V_{\pi}(s) = \mathbb{E}_{\pi}\!\left[G_t \mid S_t=s\right].
    \label{eq:value_function}
\end{equation}

\noindent\textbf{Donde:}
\begin{itemize}
    \item $V_{\pi}(s)$ representa el valor esperado del estado $s$ bajo la política $\pi$.
    \item $G_t$ representa el retorno acumulado desde el instante $t$.
    \item $S_t=s$ indica que el agente se encuentra en el estado $s$ en el instante $t$.
    \item $\mathbb{E}_{\pi}$ representa la esperanza considerando acciones tomadas según la política $\pi$.
\end{itemize}

Conceptualmente, $V_{\pi}(s)$ responde a la pregunta: ``¿qué tan bueno es estar en este
estado si sigo actuando con esta política?''. En DARL, un estado con drift severo y bajo
rendimiento debería tener menor valor si no se toman acciones correctivas adecuadas.

La función de valor de acción cuantifica la recompensa esperada al tomar una acción $a$
en un estado $s$ y luego seguir la política $\pi$:

\begin{equation}
    Q_{\pi}(s,a) = \mathbb{E}_{\pi}\!\left[G_t \mid S_t=s, A_t=a\right].
    \label{eq:q_function}
\end{equation}

\noindent\textbf{Donde:}
\begin{itemize}
    \item $Q_{\pi}(s,a)$ representa el valor esperado de ejecutar la acción $a$ en el estado $s$ bajo la política $\pi$.
    \item $S_t=s$ indica el estado actual.
    \item $A_t=a$ indica la acción ejecutada en el estado actual.
    \item $G_t$ representa el retorno acumulado posterior.
\end{itemize}

Conceptualmente, $Q_{\pi}(s,a)$ permite comparar acciones disponibles en un mismo estado.
Para DARL, esta función expresa si conviene diferir, actualizar una etapa o reentrenar el
\textit{pipeline} completo ante una señal de drift.

La política óptima satisface la ecuación de optimalidad de Bellman:

\begin{equation}
    Q^*(s,a) = \mathbb{E}\!\left[r + \gamma \max_{a'} Q^*(s',a') \mid s,a\right].
    \label{eq:bellman_optimality}
\end{equation}

\noindent\textbf{Donde:}
\begin{itemize}
    \item $Q^*(s,a)$ representa el valor óptimo de tomar la acción $a$ en el estado $s$.
    \item $r$ representa la recompensa inmediata recibida después de ejecutar la acción.
    \item $\gamma$ representa el factor de descuento.
    \item $s'$ representa el siguiente estado alcanzado.
    \item $a'$ representa una posible acción futura.
    \item $\max_{a'} Q^*(s',a')$ representa el mejor valor alcanzable desde el siguiente estado.
\end{itemize}

Conceptualmente, la ecuación de Bellman expresa que una acción es buena si produce una
buena recompensa inmediata y deja al agente en un estado desde el cual puede obtener
buenas recompensas futuras. Esta idea es fundamental para aprender políticas de
mantenimiento que consideren efectos a largo plazo \cite{sutton2018rl}.
 
\subsection{Proceso de Decisión de Markov Parcialmente Observable (POMDP)}
\label{sec:pomdp}
 
En muchos problemas reales, el agente no puede observar el estado completo del entorno,
sino solamente una señal parcial y ruidosa. Este escenario es capturado por el Proceso
de Decisión de Markov Parcialmente Observable (POMDP), definido como:

\begin{equation}
    \mathcal{P} = \langle \mathcal{S}, \mathcal{A}, \mathcal{O}, P, R, \Omega, \gamma \rangle.
    \label{eq:pomdp_tuple}
\end{equation}
 
\noindent\textbf{Donde:}
\begin{itemize}
    \item $\mathcal{P}$ representa el POMDP.
    \item $\mathcal{S}$ representa el espacio de estados reales del entorno.
    \item $\mathcal{A}$ representa el espacio de acciones disponibles para el agente.
    \item $\mathcal{O}$ representa el espacio de observaciones disponibles.
    \item $P$ representa la función de transición entre estados.
    \item $R$ representa la función de recompensa.
    \item $\Omega$ representa la función de observación.
    \item $\gamma$ representa el factor de descuento.
\end{itemize}

Conceptualmente, un POMDP extiende el MDP para casos en los que el agente no observa
directamente el estado real. Esto encaja con DARL porque el estado real de las distribuciones
del \textit{pipeline} no es accesible completamente; el agente solo observa métricas estimadas
a partir de muestras finitas \cite{kaelbling1998pomdp}.

La función de observación se define como:

\begin{equation}
    \Omega(o \mid s',a) = \Pr(O_{t+1}=o \mid S_{t+1}=s', A_t=a).
    \label{eq:observation_function}
\end{equation}

\noindent\textbf{Donde:}
\begin{itemize}
    \item $\Omega(o \mid s',a)$ representa la probabilidad de recibir la observación $o$ después de llegar al estado $s'$ tras ejecutar la acción $a$.
    \item $O_{t+1}$ representa la observación recibida en el siguiente instante.
    \item $S_{t+1}$ representa el estado real alcanzado.
    \item $A_t$ representa la acción ejecutada por el agente.
\end{itemize}

Conceptualmente, esta fórmula representa la diferencia entre el estado real y lo que el
agente puede medir. En DARL, las métricas PSI, KS, C2ST y severidad son observaciones
imperfectas del verdadero estado de drift y degradación del \textit{pipeline}.

Dado que el estado real no es directamente accesible, el agente puede mantener una
distribución de probabilidad sobre los estados posibles, denominada \textit{belief state}:

\begin{equation}
    b_t(s) = P(s_t=s \mid o_1,a_1,\ldots,o_t).
    \label{eq:belief_state}
\end{equation}
 
\noindent\textbf{Donde:}
\begin{itemize}
    \item $b_t(s)$ representa la creencia del agente sobre la probabilidad de estar en el estado $s$ en el instante $t$.
    \item $s_t$ representa el estado real en el instante $t$.
    \item $o_1,\ldots,o_t$ representan las observaciones recibidas hasta el instante $t$.
    \item $a_1,\ldots,a_{t-1}$ representan las acciones ejecutadas previamente.
    \item $P(\cdot \mid \cdot)$ representa una probabilidad condicional.
\end{itemize}

Conceptualmente, el \textit{belief state} resume la historia de observaciones y acciones
en una estimación probabilística del estado actual. En DARL, esto justifica que el agente
aprenda a decidir bajo incertidumbre, ya que las métricas observadas pueden ser ruidosas
o insuficientes para determinar con certeza el tipo de drift.

La política óptima en un POMDP puede expresarse como una función del \textit{belief state}:

\begin{equation}
    \pi^*(b_t).
    \label{eq:pomdp_policy}
\end{equation}

\noindent\textbf{Donde:}
\begin{itemize}
    \item $\pi^*$ representa la política óptima.
    \item $b_t$ representa el \textit{belief state} en el instante $t$.
\end{itemize}

Conceptualmente, esta notación indica que el agente no decide a partir del estado real,
sino a partir de su mejor estimación del estado. En implementaciones prácticas como DARL,
esta estimación se aproxima mediante el vector de métricas observado.
 
En el contexto de DARL, el estado oculto $s_t$ es el estado real de las distribuciones del
\textit{pipeline}, incluyendo la distribución real $P_t(X)$, la relación real $P_t(Y \mid X)$
y el grado de degradación del modelo. Lo que sí se observa es el vector de métricas
estadísticas calculadas sobre muestras finitas:
 
\begin{equation}
    \mathbf{o}_t = \left[
    \overline{\text{PSI}}_t,\;
    \overline{D}_{\text{KS},t},\;
    \text{C2ST}_t,\;
    S_t
    \right] \in \mathbb{R}^{m}.
    \label{eq:darl_observation_vector}
\end{equation}
 
\noindent\textbf{Donde:}
\begin{itemize}
    \item $\mathbf{o}_t$ representa la observación disponible para DARL en el instante $t$.
    \item $\overline{\text{PSI}}_t$ representa el promedio de PSI sobre las variables monitoreadas.
    \item $\overline{D}_{\text{KS},t}$ representa el promedio de los estadísticos KS sobre las variables monitoreadas.
    \item $\text{C2ST}_t$ representa la señal global del test de dos muestras basado en clasificación.
    \item $S_t$ representa el puntaje combinado de severidad del drift.
    \item $\mathbb{R}^{m}$ representa el espacio vectorial de observaciones de dimensión $m$.
\end{itemize}

Conceptualmente, esta fórmula muestra cómo DARL convierte señales estadísticas de
monitoreo en una observación compacta para el agente. Debido a que estas métricas se
calculan sobre ventanas finitas, contienen ruido e incertidumbre; por ello, el problema se
modela de forma natural como un POMDP.
 
% ------------------------------------------------------------
\section{Proximal Policy Optimization (PPO)}
\label{sec:ppo}
% ------------------------------------------------------------
 
PPO es un algoritmo de aprendizaje por refuerzo de la familia de métodos de
\textit{policy gradient} que optimiza directamente la política $\pi_\theta$ con respecto
al objetivo de la Ecuación~\eqref{eq:objetivo_rl} \cite{schulman2017ppo}.
 
\subsection{Métodos de policy gradient}
 
Los métodos de \textit{policy gradient} parametrizan la política como $\pi_\theta$ y
estiman el gradiente del objetivo mediante:
 
\begin{equation}
    \nabla_\theta J(\theta)
    = \mathbb{E}_{\pi_\theta}\!\left[
    \nabla_\theta \log \pi_\theta(a_t \mid s_t)\,\hat{A}_t
    \right].
    \label{eq:policy_gradient}
\end{equation}
 
\noindent\textbf{Donde:}
\begin{itemize}
    \item $J(\theta)$ representa el objetivo de desempeño de la política parametrizada.
    \item $\theta$ representa los parámetros de la política.
    \item $\nabla_\theta J(\theta)$ representa el gradiente del objetivo respecto a los parámetros.
    \item $\pi_\theta(a_t \mid s_t)$ representa la probabilidad de elegir la acción $a_t$ en el estado $s_t$ bajo la política actual.
    \item $\nabla_\theta \log \pi_\theta(a_t \mid s_t)$ indica cómo cambiar los parámetros para aumentar o reducir la probabilidad de la acción tomada.
    \item $\hat{A}_t$ representa la ventaja estimada de la acción tomada en el instante $t$.
    \item $\mathbb{E}_{\pi_\theta}$ representa la esperanza sobre trayectorias generadas por la política $\pi_\theta$.
\end{itemize}

Conceptualmente, esta fórmula indica que una acción debe volverse más probable si produjo
una ventaja positiva y menos probable si produjo una ventaja negativa. El gradiente no solo
considera qué acción se tomó, sino si esa acción fue mejor o peor que lo esperado en ese
estado \cite{sutton2018rl}.

La ventaja estimada puede expresarse como:

\begin{equation}
    \hat{A}_t = Q^{\pi_\theta}(s_t,a_t) - V^{\pi_\theta}(s_t).
    \label{eq:advantage}
\end{equation}

\noindent\textbf{Donde:}
\begin{itemize}
    \item $\hat{A}_t$ representa la ventaja estimada en el instante $t$.
    \item $Q^{\pi_\theta}(s_t,a_t)$ representa el valor esperado de tomar la acción $a_t$ en el estado $s_t$ y luego seguir la política $\pi_\theta$.
    \item $V^{\pi_\theta}(s_t)$ representa el valor esperado del estado $s_t$ bajo la política $\pi_\theta$.
\end{itemize}

Conceptualmente, la ventaja mide cuánto mejor o peor fue una acción respecto al promedio
esperado en ese estado. En DARL, si reentrenar todo el \textit{pipeline} produce alta
recuperación pero con costo excesivo cuando bastaba actualizar $f_\phi$, su ventaja puede
ser menor que la de una acción selectiva más eficiente.
 
\subsection{La función objetivo de PPO}
 
PPO resuelve el problema de inestabilidad imponiendo una restricción sobre el tamaño
del cambio de política mediante el \textit{clipping} del ratio de probabilidad
\cite{schulman2017ppo}. El ratio entre la política nueva y la política anterior se define como:

\begin{equation}
    r_t(\theta) =
    \frac{\pi_\theta(a_t \mid s_t)}
    {\pi_{\theta_{\text{old}}}(a_t \mid s_t)}.
    \label{eq:ppo_ratio}
\end{equation}

\noindent\textbf{Donde:}
\begin{itemize}
    \item $r_t(\theta)$ representa el cambio relativo en la probabilidad de la acción $a_t$ bajo la política nueva respecto a la política anterior.
    \item $\pi_\theta(a_t \mid s_t)$ representa la probabilidad asignada por la política actual.
    \item $\pi_{\theta_{\text{old}}}(a_t \mid s_t)$ representa la probabilidad asignada por la política antes de la actualización.
    \item $\theta_{\text{old}}$ representa los parámetros de la política anterior.
\end{itemize}

Conceptualmente, este ratio mide cuánto está cambiando la política para la acción que fue
ejecutada. Si el ratio se aleja demasiado de $1$, significa que la actualización está modificando
de forma agresiva el comportamiento del agente.

La función objetivo recortada de PPO es:
 
\begin{equation}
    L^{\text{CLIP}}(\theta)
    = \mathbb{E}_t\!\left[
    \min\!\left(
    r_t(\theta)\hat{A}_t,
    \text{clip}\!\left(r_t(\theta),1-\epsilon,1+\epsilon\right)\hat{A}_t
    \right)
    \right].
    \label{eq:ppo_clip}
\end{equation}
 
\noindent\textbf{Donde:}
\begin{itemize}
    \item $L^{\text{CLIP}}(\theta)$ representa la función objetivo recortada que PPO maximiza.
    \item $\mathbb{E}_t$ representa el promedio empírico sobre los pasos de tiempo del \textit{rollout}.
    \item $r_t(\theta)$ representa el ratio de probabilidad entre la política nueva y la política anterior.
    \item $\hat{A}_t$ representa la ventaja estimada.
    \item $\epsilon \in (0,1)$ es un hiperparámetro que controla el tamaño máximo permitido del cambio de política.
    \item $\text{clip}(x,a,b)$ restringe el valor $x$ al intervalo $[a,b]$.
    \item $\min(\cdot)$ selecciona el término más conservador entre la actualización sin recorte y la actualización recortada.
\end{itemize}

Conceptualmente, PPO evita que la política cambie demasiado en una sola actualización.
Cuando la ventaja es positiva, el agente no puede aumentar indefinidamente la probabilidad
de la acción; cuando la ventaja es negativa, tampoco puede reducirla de forma excesiva. Esta
cota pesimista mejora la estabilidad del entrenamiento, lo cual es importante en DARL porque
cada episodio puede implicar costos de reentrenamiento no triviales.
 
\subsection{Arquitectura Actor-Crítico}
 
PPO implementa una arquitectura Actor-Crítico \cite{sutton2018rl,schulman2017ppo}. En
DARL, sus dos componentes principales pueden expresarse como:

\begin{equation}
    \pi_\theta(a \mid \mathbf{o}_t), \qquad V_\psi(\mathbf{o}_t).
    \label{eq:actor_critic}
\end{equation}

\noindent\textbf{Donde:}
\begin{itemize}
    \item $\pi_\theta(a \mid \mathbf{o}_t)$ representa el Actor, es decir, la política que asigna probabilidades a las acciones posibles dado el vector de observación.
    \item $\theta$ representa los parámetros del Actor.
    \item $V_\psi(\mathbf{o}_t)$ representa el Crítico, es decir, la estimación del valor de la observación actual.
    \item $\psi$ representa los parámetros del Crítico.
    \item $\mathbf{o}_t$ representa el vector de métricas observado por DARL en el instante $t$.
\end{itemize}

Conceptualmente, el Actor decide qué acción tomar y el Crítico evalúa qué tan buena es la
situación observada. Esta separación permite estimar ventajas con menor varianza y mejorar
la estabilidad del aprendizaje. En DARL, el Actor produce una distribución sobre las acciones
posibles: diferir, actualizar \textit{features}, actualizar modelo o reentrenar todo el \textit{pipeline}.
 
\subsection{Justificación de PPO para DARL}
 
PPO es el algoritmo elegido para DARL por tres razones específicas al problema
\cite{schulman2017ppo,sutton2018rl}:
 
\begin{enumerate}
    \item Espacio de acciones discreto y pequeño. El espacio de acciones de DARL tiene cardinalidad:

    \begin{equation}
        |\mathcal{A}| = 4.
        \label{eq:action_space_cardinality}
    \end{equation}

    \noindent\textbf{Donde:}
    \begin{itemize}
        \item $|\mathcal{A}|$ representa el número de acciones disponibles.
        \item $4$ corresponde a las acciones: diferir, actualizar preprocesamiento, actualizar modelo y reentrenar ambas etapas.
    \end{itemize}

    Conceptualmente, un espacio de acciones pequeño se ajusta naturalmente a una política categórica implementada con una capa \textit{softmax}.
 
    \item Eficiencia muestral. Cada episodio de entrenamiento puede implicar reentrenar XGBoost, lo que tiene un costo computacional no trivial. PPO es un algoritmo \textit{on-policy} que aprovecha cada \textit{rollout} mediante múltiples épocas de actualización sobre los datos recolectados.
 
    \item Estabilidad de entrenamiento. El mecanismo de \textit{clipping} de la Ecuación~\eqref{eq:ppo_clip} previene cambios excesivos de política durante el entrenamiento, lo cual es crítico cuando el número de episodios disponibles es limitado.
\end{enumerate}
 
% ------------------------------------------------------------
\section{XGBoost como modelo predictivo base}
\label{sec:xgboost}
% ------------------------------------------------------------
 
XGBoost (\textit{eXtreme Gradient Boosting}) es un algoritmo de aprendizaje basado en
árboles de decisión potenciados por gradiente que ha demostrado ser altamente competitivo
en tareas de clasificación y regresión sobre datos tabulares
\cite{chen2016xgboost,gardner2023tableshift}.
 
El modelo construye un ensamble aditivo de $K$ árboles de decisión:
 
\begin{equation}
    \hat{y}_i = \sum_{k=1}^{K} f_k(x_i),
    \qquad f_k \in \mathcal{F}.
    \label{eq:xgboost_prediction}
\end{equation}
 
\noindent\textbf{Donde:}
\begin{itemize}
    \item $\hat{y}_i$ representa la predicción del modelo para la observación $i$.
    \item $x_i$ representa la observación de entrada $i$.
    \item $K$ representa el número total de árboles del ensamble.
    \item $f_k$ representa el árbol de decisión número $k$.
    \item $\mathcal{F}$ representa el espacio de árboles de regresión disponibles para el modelo.
\end{itemize}

Conceptualmente, XGBoost no aprende un único árbol, sino una suma de árboles que corrigen
progresivamente los errores de los árboles anteriores. Esta estructura lo hace flexible para
capturar relaciones no lineales y combinaciones complejas de variables en datos tabulares.
 
El entrenamiento minimiza una función objetivo regularizada:
 
\begin{equation}
    \mathcal{L}(\theta) =
    \sum_{i=1}^{n} \ell(y_i,\hat{y}_i)
    + \sum_{k=1}^{K} \Omega(f_k).
    \label{eq:xgboost_objective}
\end{equation}
 
\noindent\textbf{Donde:}
\begin{itemize}
    \item $\mathcal{L}(\theta)$ representa la función objetivo que XGBoost busca minimizar.
    \item $\theta$ representa los parámetros del modelo.
    \item $n$ representa el número de observaciones de entrenamiento.
    \item $y_i$ representa la etiqueta real de la observación $i$.
    \item $\hat{y}_i$ representa la predicción del modelo para la observación $i$.
    \item $\ell(y_i,\hat{y}_i)$ representa la pérdida de predicción, por ejemplo, entropía cruzada en clasificación binaria.
    \item $\Omega(f_k)$ representa la penalización de complejidad asociada al árbol $f_k$.
\end{itemize}

Conceptualmente, esta función objetivo equilibra dos metas: ajustar bien los datos y evitar
modelos excesivamente complejos. La primera suma mide el error de predicción; la segunda
suma penaliza árboles demasiado grandes o con pesos muy altos, reduciendo el riesgo de
sobreajuste \cite{chen2016xgboost}.

La penalización de complejidad de cada árbol se define como:

\begin{equation}
    \Omega(f_k) = \gamma_{\text{xgb}} T_k
    + \frac{1}{2}\lambda_{\text{xgb}}\lVert\mathbf{w}_k\rVert^2.
    \label{eq:xgboost_regularization}
\end{equation}

\noindent\textbf{Donde:}
\begin{itemize}
    \item $\Omega(f_k)$ representa la complejidad penalizada del árbol $k$.
    \item $\gamma_{\text{xgb}}$ representa el costo asociado a añadir hojas al árbol.
    \item $T_k$ representa el número de hojas del árbol $k$.
    \item $\lambda_{\text{xgb}}$ representa el peso de regularización sobre los valores de las hojas.
    \item $\mathbf{w}_k$ representa el vector de pesos o valores asignados a las hojas del árbol $k$.
    \item $\lVert\mathbf{w}_k\rVert^2$ representa la norma cuadrada de dichos pesos.
\end{itemize}

Conceptualmente, esta regularización penaliza tanto la estructura del árbol como la magnitud
de sus predicciones en las hojas. Esto permite que XGBoost mantenga capacidad predictiva sin
crecer de forma innecesaria. En DARL, esta propiedad es útil porque el costo de actualizar o
reentrenar el modelo puede medirse y compararse contra la recuperación de rendimiento.
 
En el contexto del \textit{pipeline} de dos etapas, los parámetros relevantes son:

\begin{equation}
    \{\phi,\theta\}.
    \label{eq:pipeline_parameters}
\end{equation}

\noindent\textbf{Donde:}
\begin{itemize}
    \item $\phi$ representa los parámetros de la etapa de preprocesamiento, como medias, desviaciones estándar, cuantiles, imputaciones y codificaciones.
    \item $\theta$ representa los parámetros de la etapa predictiva, en este caso asociados al modelo XGBoost.
\end{itemize}

Conceptualmente, esta notación resume la idea central de la tesis: el mantenimiento del
\textit{pipeline} no tiene por qué actualizar todos los parámetros siempre. Cada acción de
DARL modifica un subconjunto distinto de $\{\phi,\theta\}$ con los datos de la ventana actual.
Así, el agente puede aprender cuándo basta actualizar el preprocesamiento, cuándo es necesario
actualizar el modelo y cuándo conviene reentrenar ambas etapas.
